% Supplementary material of the LTCC manuscript (Aerospace Science and Technology submission).
% Build after main.tex: pdflatex supplementary && bibtex supplementary && pdflatex supplementary (x2)
% Every number comes from the generated tables/numbers*.tex files; none is typed by hand.
\documentclass[a4paper,10pt]{article}
\usepackage[margin=2.2cm]{geometry}
\usepackage[T1]{fontenc}
\usepackage{amsmath}
\usepackage{newtxtext,newtxmath}
\usepackage[numbers]{natbib}
\usepackage{graphicx}
\usepackage{booktabs}
\usepackage{threeparttable}
\usepackage{array}
\usepackage{xcolor}
\usepackage{xr-hyper}
\usepackage[hidelinks]{hyperref}
\externaldocument{main}
% propositions and proofs in upright type
\newcounter{prop}
\renewcommand{\theprop}{S\arabic{prop}}
\newenvironment{prop}{\par\smallskip\refstepcounter{prop}\noindent Proposition~\theprop.\ }{\par\smallskip}
\providecommand{\QED}{\hfill$\blacksquare$}
\newenvironment{proof}{\par\noindent Proof.\ }{\QED\par}
\newcommand{\todo}[1]{\textcolor{red}{[TODO: #1]}}
\newcommand{\mname}{LTCC}  % name of the proposed method
\newcommand{\val}[1]{\ifcsname v#1\endcsname\csname v#1\endcsname\else\todo{#1}\fi}
\makeatletter
\@for\@numfile:=static_numbers,static_numbers_XL,static_numbers_scale,numbers,numbers_s2,numbers_s3,numbers_s3x,%
numbers_err,numbers_s4,numbers_s4c,numbers_s4x,numbers_s4cx,numbers_latepen,numbers_s5,numbers_s6,numbers_extra,%
numbers_demo,numbers_robust,numbers_calib,numbers_s7,numbers_s7x,numbers_s7c\do{%
  \IfFileExists{tables/\@numfile.tex}{\input{tables/\@numfile.tex}}{}}
\makeatother
\newsavebox{\ltccfigbox}
\newcommand{\fig}[2]{\sbox{\ltccfigbox}{\includegraphics{figures/#1.pdf}}%
  \ifdim\wd\ltccfigbox>\linewidth\resizebox{\linewidth}{!}{\usebox{\ltccfigbox}}\else\usebox{\ltccfigbox}\fi}
\newcommand{\uprighttext}{\def\textbf##1{##1}\def\textit##1{##1}}
\newcommand{\inputtable}[1]{\begingroup\uprighttext\input{tables/#1.tex}\endgroup}
\renewcommand{\thesection}{S\arabic{section}}
\renewcommand{\thesubsection}{\thesection.\arabic{subsection}}
\renewcommand{\thetable}{S\arabic{table}}
\renewcommand{\thefigure}{S\arabic{figure}}
\renewcommand{\theequation}{S\arabic{equation}}

\begin{document}
\title{Supplementary material for: Chance Constraints as Effective Deadlines for Fast Replanning of UAV Fleets Under Uncertainty}
\author{Cao-Phuc Ha, Vu Minh Anh Nguyen, Duong Thanh Tran, and Phu-Nguyen Le}
\date{}
\maketitle

\begin{abstract}
This document contains the parts of the paper that the main text only summarises: the planning deadlines and the proofs
of Propositions 1 and 2, the stopping rule and its two delay costs, the scenario generator and the fitted error models,
the parameters of the compared methods and further details of the protocol, the exact-execution benchmark of the
stopping rules, the full results of study~1, further results of study~2 per regime and in tier C, the regimes of
study~3, the reliability of the error model in study~4, the re-scoring that motivated study~5, the calibration of the
predicted probability in studies 4-7, the sensitivity of all tier B tests to the resampling unit and further results
of studies 2-7. Section, equation and table numbers without the prefix S refer to the main paper.
\end{abstract}

\section{Planning deadlines and proofs}\label{sec:Sproofs}
This section states the planning deadlines of Section~\ref{sec:method} of the main paper in full and proves
Propositions~\ref{prop:cc} and~\ref{prop:costfb}. For task $j$, $f_j$ is the planned finish time, $d_j$ the deadline,
$w_j$ the reward, $h_j=f_j-t_0$ the lead time and $F_j=f_j+e_j$ the actual finish time, with
$e_j\sim\mathcal{N}(m(h_j),s(h_j)^2)$ and $m$, $s$ as in Eq.~\eqref{eq:errmodel}. The effective deadline of
Eq.~\eqref{eq:cc} at safety factor $z$ is written $\tilde d_j(z)$, and $\tilde d_j(-\infty)$ is read as $d_j$.

\subsection{Planning deadlines}
Let $\tilde H$ be the set of tasks that no vehicle can finish by $\tilde d_j(z)$ on a direct flight, and $H$ the set of
tasks that no vehicle can finish by $d_j$. The nominal fallback plans against
\begin{equation}\label{eq:Sfallback}
d'_j=\begin{cases} d_j, & j\in\tilde H\setminus H,\\ \tilde d_j(z), & \text{otherwise.}\end{cases}
\end{equation}
For the cost-aware fallback, let $\theta_j=c_j/(w_j+c_j)$ be the break-even probability of Eq.~\eqref{eq:theta},
$z_j=\Phi^{-1}(\theta_j)$ with $z_j=-\infty$ for $c_j=0$, and $z_{p,j}=\max(z,z_j)$. With $\tilde H_p$ the set of tasks
that no vehicle can finish by $\tilde d_j(z_{p,j})$ on a direct flight, the planning deadline is
\begin{equation}\label{eq:Scostfb}
d'_j=\begin{cases} \min\{d_j,\tilde d_j(z_j)\}, & j\in\tilde H_p\setminus H,\\ \tilde d_j(z_{p,j}), &
\text{otherwise.}\end{cases}
\end{equation}
A task that cannot meet this deadline either is given up. Finding $\tilde H_p$ and $H$ requires one direct-flight check
per task and vehicle. Where $1+b_1+z_js_1\le0.05$, the effective deadline is not defined and the nominal deadline is
used, without the guarantee of part (iii) of Proposition~\ref{prop:costfb}. With the error models of this paper, this
happens only for $\theta_j<0.03$, which none of the studied late costs produces. Every planner caps the effective deadline
at $d_j$, so that no task is planned beyond its nominal deadline.

\subsection{Proposition 1}
In full, Proposition~\ref{prop:cc} states the following. Let $1+b_1+zs_1>0$ and $s_0,s_1\ge0$. (i) A finish time
satisfies $f_j+m(f_j-t_0)+z\,s(f_j-t_0)\le d_j$ if and only if $f_j\le\tilde d_j(z)$. (ii) Under the error model, a task
planned to finish at $f_j\le\tilde d_j(z)$ meets its deadline with probability at least $\Phi(z)$ if its route starts as
planned. (iii) With $b_0=b_1=s_0=s_1=0$, $\tilde d_j(z)=d_j$ and the planning model is the nominal one.

\begin{proof}
Substituting the error model gives $f_j(1+b_1+zs_1)\le d_j-b_0-zs_0+(b_1+zs_1)t_0$, and dividing by the positive factor
gives (i). For (ii), $F_j=f_j+e_j$ with $e_j\sim\mathcal{N}(m(h_j),s(h_j)^2)$. By (i), $d_j-f_j-m(h_j)\ge z\,s(h_j)$. If
$s(h_j)>0$, then $\Pr[F_j\le d_j]=\Phi\big((d_j-f_j-m(h_j))/s(h_j)\big)\ge\Phi(z)$. If $s(h_j)=0$, which happens only
for $s_0=0$ and $h_j=0$, the error equals $m(h_j)$ with certainty, $F_j\le d_j$ holds, and the probability is one. Part
(iii) is immediate.
\end{proof}

The cap at $d_j$ matters only where $\tilde d_j(z)>d_j$, which happens when $z<0$ or the mean error is negative. It then
adds the nominal deadline to the constraint, which becomes tighter, so part (ii) still holds, whereas the equivalence of
part (i) refers to the uncapped deadline.

\subsection{Proposition 2}
In full, Proposition~\ref{prop:costfb} states the following. Let $s_0,s_1\ge0$, $1+b_1+zs_1>0$ and $1+b_1+z_js_1>0$
for every finite $z_j$. (i) If $c_j=0$ for every task, Eq.~\eqref{eq:Scostfb} is the nominal fallback of
Eq.~\eqref{eq:Sfallback}. (ii) A task with $\theta_j\ge\Phi(z)$ never falls back, and it is planned against
$\tilde d_j(z_j)$. (iii) Under the error model, every task planned to finish by $d'_j$ has an expected contribution
$w_j\Pr[F_j\le d_j]-c_j\Pr[F_j>d_j]\ge0$ if its route starts as planned. (iv) If every $c_j$ is non-decreasing in
$\lambda$, every planning deadline $d'_j$ is non-increasing in $\lambda$.

\begin{proof}
(i) $z_j=-\infty$ gives $z_{p,j}=z$ and $\min\{d_j,\tilde d_j(-\infty)\}=d_j$. (ii) $z_j\ge z$ holds if and only if
$\theta_j\ge\Phi(z)$. Then $z_{p,j}=z_j$, and a task that cannot meet $\tilde d_j(z_{p,j})$ cannot meet
$\min\{d_j,\tilde d_j(z_j)\}$ either, so it is given up. (iii) A protected task meets its deadline with probability at
least $\Phi(z_{p,j})\ge\Phi(z_j)$, and a fallback task with probability at least $\Phi(z_j)$, both by
Proposition~\ref{prop:cc}, whose part (ii) also holds with the cap at $d_j$. Hence $p\ge\theta_j$, and the expected
contribution $(w_j+c_j)p-c_j$ is not negative. (iv) $\theta_j$ and $z_j$ do not decrease with $\lambda$, and with
$s_0,s_1\ge0$, $\tilde d_j(z)$ does not increase with $z$. A fallback task stays one when $\lambda$ grows, because
$z_{p,j}$ does not decrease. A protected task can become a fallback task only while $z_j<z$, and then $z_{p,j}=z$ does not
change. In each case the deadline does not increase.
\end{proof}

Part (iii) is a local criterion. It prices each attempt on its own and ignores the delay that a fallback task causes to
the tasks after it on the same route, as well as the reassignments it may trigger, so it does not guarantee that
attempting the task raises the utility of the whole plan.


\section{Stopping rule and delay costs}\label{sec:Sstop}
\subsection{Learned stopping rule}
The objective of a plan $\pi$ in Section~\ref{sec:problem} of the main paper is $J(\pi)=\alpha R(\pi)-\gamma
E(\pi)-\delta C(\pi)$, where $R$ is the reward of the tasks served on time, $E$ the flown distance in kilometres and $C$
the number of tasks that move to another UAV, with $\alpha=1$, $\gamma=0.05$ per kilometre and $\delta=0.3$. $W$ denotes
the total reward of the open tasks. Every 5\,ms, the stopping rule compares the predicted gain of searching for another $H=10$\,ms with the cost of doing so.
With $\widehat{\Delta J}$ as the predicted gain, the search stops at time $\tau$ if
\begin{equation}\label{eq:Sstop}
\widehat{\Delta J}\;\le\;C_{\mathrm{delay}}(\tau)+\mu W H .
\end{equation}
The gain is predicted from the search history by a Poisson generalised linear model with a log
link~\cite{mccullagh1989glm} and a small $L_2$ penalty. Its eleven features describe the elapsed time, the number of
iterations, the recent improvements over 5\,ms, over 20\,ms and over the second half of the elapsed time, the time since
the last improvement, the total improvement over the repaired plan, the unassigned share of the reward, the numbers of
tasks and UAVs, and the iteration rate. The target is the improvement in the next 10\,ms in percent of $W$. The model is
fitted on 500\,ms search trajectories of development scenarios only, and a prediction costs one dot product. The search
loop enforces the hard time budget $B$ independently of the stopping rule. The price of computation is
$\mu=10^{-5}$ per millisecond.

\subsection{Deterministic delay cost}
The delay cost is computed from the best plan found so far. A late start of a route is partly absorbed by the waiting
times along it, in the same way as the forward time slack of routing with time windows~\cite{savelsbergh1992vrptw}. The
tolerance $\mathrm{tol}_j$ of a task is the margin between its finish time and its deadline, plus all waiting times on
the route up to this task. A task is lost when the start shift $\sigma_j(\tau)$ of its route exceeds its tolerance. The
start shift is zero for a UAV that is still busy, and it equals the waiting time until the publication for an idle UAV.
With $v_j=\alpha w_j-\delta\,\mathbb{1}[\text{$j$ changes UAV}]$ as the value of task $j$ in the objective and
$\Lambda(\tau)=\sum_j v_j\,\mathbb{1}[\sigma_j(\tau)>\mathrm{tol}_j]$, the delay cost of searching $H$ milliseconds longer
is $C_{\mathrm{delay}}(\tau)=\Lambda(\tau+H)-\Lambda(\tau)$.

\begin{prop}\label{prop:delay}
For a feasible plan executed exactly, (i) $\Lambda(\tau)$ is the value lost through missed deadlines if every UAV
follows its route, (ii) the loss of the objective is at most $\Lambda(\tau)$ if late tasks are skipped, and
(iii) $\Lambda(\tau)$ does not decrease with $\tau$ and is zero while all UAVs are busy with their committed tasks.
Hence $C_{\mathrm{delay}}$ is non-negative and vanishes while no UAV is waiting.
\end{prop}
\begin{proof}
A route that starts later by some shift reaches its first task later by the same amount. Each waiting time absorbs part
of the delay, so the delay at the end of a task equals the start shift minus all earlier waiting times, if this
difference is positive. The task therefore misses its deadline exactly when the start shift exceeds its tolerance.
Waiting consumes no energy, which proves (i). If late tasks are skipped, the triangle inequality ensures that every
later task is reached no later and with no more energy. All tasks that are on time in (i) remain on time, and the flown
distance, hence the energy term $\gamma E$, does not grow. The objective therefore loses the values $v_j$ of the skipped
tasks and nothing else, which is at most $\Lambda(\tau)$ and proves (ii). Part (iii) holds because the start shift does
not decrease with $\tau$ and every tolerance is non-negative in a feasible plan.
\end{proof}

\subsection{Execution-aware delay cost}\label{sec:Srisk}
Under uncertain execution, task $j$ is lost when $\sigma_j(\tau)+e_j>\mathrm{tol}_j$, where $e_j$ is its finish-time
error. The execution-aware cost replaces the indicator in $\Lambda$ by the probability of this event. The error of task
$j$ is modelled as $e_j\sim\mathcal{N}(m(h_j),s(h_j)^2)$, where $h_j=f_j-t_0$ is its lead time, the nominal finish time
minus the time of the planning snapshot. The mean and the spread are linear in the lead time,
\begin{equation}\label{eq:Serrmodel}
m(h)=b_0+b_1h,\qquad s(h)=s_0+s_1h,
\end{equation}
because errors accumulate along a route. The expected dissemination latency $\bar L$ is added to the start shift. The
expected loss and the execution-aware delay cost are
\begin{equation}\label{eq:Srisk}
\Lambda_{\mathrm{R}}(\tau)=\sum_j v_j\,\Phi\!\left(\frac{\sigma_j(\tau)+\bar L+m(h_j)-\mathrm{tol}_j}{s(h_j)}\right),
\qquad C_{\mathrm{R}}(\tau)=\Lambda_{\mathrm{R}}(\tau+H)-\Lambda_{\mathrm{R}}(\tau),
\end{equation}
where $\Phi$ is the standard normal distribution function. Search and predictor are unchanged, and $C_{\mathrm{R}}$
replaces $C_{\mathrm{delay}}$ in Eq.~\eqref{eq:Sstop}.

\begin{prop}\label{prop:risk}
(i) With $b_0=b_1=s_0=s_1=0$ and $\bar L=0$, $\Lambda_{\mathrm{R}}=\Lambda$, and the stopping rule takes exactly the
decisions of the deterministic one. (ii) $\Lambda_{\mathrm{R}}(\tau)$ does not decrease with $\tau$, so $C_{\mathrm{R}}\ge0$.
(iii) Under the error model, $\Lambda_{\mathrm{R}}(\tau)$ is the expected value of the tasks whose deadline is missed if
every UAV follows its route.
\end{prop}
\begin{proof}
With zero spread, $\Phi$ is read as the indicator of a positive argument, which gives (i). The start shift does not
decrease with $\tau$ and $\Phi$ is non-decreasing, which gives (ii). Part (iii) follows from the proof of
Proposition~\ref{prop:delay}, because task $j$ misses its deadline exactly when its delay at the finish,
$\sigma_j+e_j$ minus the absorbed waiting time, is positive, and the expectation of the indicator of this event is the
probability in Eq.~\eqref{eq:Srisk}.
\end{proof}

The parameters of Eq.~\eqref{eq:Serrmodel} are fitted on development missions. The mean is estimated by least squares,
and the spread by least squares on the absolute residuals scaled by $\sqrt{\pi/2}$, which is less sensitive to the heavy
tail of the errors than squared residuals. The latency $\bar L$ is the mean of the latency distribution in tier B and
the measured mean plan-acceptance latency in tier C.



\section{Scenario generator, error models and exact-execution snapshots}\label{sec:Sscen}
\subsection{Scenarios}
Missions are created by a seeded generator with the parameters in Table~\ref{tab:params}. The map is a square with a
side of 10\,km and the depot at its centre, and half of the tasks are clustered around three hotspots. Half of the tasks
are known at the start, and the others arrive in batches. Three fleet scales with 5, 10 and 20 UAVs and 25, 50 and 140
tasks are combined with three event profiles. The arrivals profile contains six arrival batches, the failures profile
three batches and two UAV failures, and the deadlines profile three batches and four tightened deadlines. A fourth scale
with 40 UAVs and 280 tasks is used for scalability only.

\begin{table}[t]
\caption{Scenario generator parameters (distances in m, times in s, energy in metres of flight).}\label{tab:params}
\centering\small
\setlength{\tabcolsep}{4pt}
\begin{tabular}{ll@{\qquad}ll}
\toprule
Map side / horizon & 10\,000 / 1\,800 & Service time & $U[30,90]$ \\
UAV speed & $U[10,16]$\,m/s & Time-window slack & $U[180,600]$ \\
Energy capacity & $U[18,28]$\,km & Release delay & $U[0,120]$ \\
Energy reserve & 10\% of capacity & Rewards 1, 2, 3, 5 & probabilities .4, .3, .2, .1 \\
Hover cost & 6 per second & Hotspots & 3, std.\ 600, 50\% of tasks \\
Tasks known at start & 50\% & Scales (UAVs/tasks) & 5/25, 10/50, 20/140, 40/280 \\
\bottomrule
\end{tabular}
\end{table}

\subsection{Static snapshots with exact execution}
The stopping rules are first compared without execution uncertainty. Each mission is simulated once with a
deterministic reference planner, and every replanning call after the start is stored with its previous plan. All
methods are evaluated on these snapshots, and regret is measured against the best known objective of each snapshot,
which includes a 3\,s ALNS run and, for the small and medium scales, the OR-Tools routing solver~\cite{ortools} with
guided local search~\cite{voudouris1999gls}. Because all ALNS variants follow the same search trajectory, whole families
of each rule are evaluated exactly offline, with 12 fixed budgets, 10 stagnation patience values and 9 compute prices,
and each baseline is compared with AnytimeSwarm at the same mean computation. This benchmark uses its own planned
test seeds 2000-2029 (\val{NsnapTest} snapshots of 270 missions), and it was reproduced on the final software platform
before the execution-aware study (Section~\ref{sec:repro}).



\inputtable{tabm_errmodels}

\section{Methods, protocol and further results}\label{sec:Sdetails}
% Details moved from the main text (methods, protocol and secondary results); every number is a \val macro.
\subsection{Error structure}
The comparison of spread models in Section~\ref{sec:method} of the main paper uses \val{ENPairsMid} and
\val{ENPairsHigh} pairs of lead time and error from \val{ENMissions} development missions per disturbance level. A
hybrid model $s(h)=\sqrt{q_0+q_1h+q_2h^2}$ also fitted worse than the linear spread, by \val{EDAicHybMid} and
\val{EDAicHybHigh} points of the Akaike information criterion. The flight-time error of a single leg had a coefficient of
variation of \val{EMpCvMid} and \val{EMpCvHigh} under mid and high disturbances, and missions explained only
\val{EShareMissionMid}\% and \val{EShareMissionHigh}\% of the residual variance.

\subsection{Parameters of the compared methods}\label{sec:Sparams}
The performance profile rule predicts the gain of further search from a static profile over the elapsed time and the
number of open tasks. The online prediction rule follows Svegliato et al.~\cite{svegliato2018meta}, but fits the
saturating curve $q(t)=a-b\,e^{-ct}$ to the best objective found so far at every check, in place of the arctangent or
logarithmic models of the original. Both use the delay cost and the price of computation of AnytimeSwarm-R. The size-scaled budget allows 1\,ms
per open task, and the stagnation rule stops after 20\% of the budget without improvement. The proportional buffer
accepts a planned finish $f_j$ if $f_j+c\,h_j\le d_j$. Moment propagation models the flight time of a leg of nominal
duration $t$ as normal with mean $(1+a)t+c_0$ and standard deviation $c_vt$, and the service time with mean $(1+a_s)s$
and standard deviation $c_ss$. Legs are independent, the start of service $\max(A_j,r_j)$ is approximated by Clark's
moments~\cite{clark1961greatest}, and a finish is accepted if $\mathrm{E}[F_j]+z\,\mathrm{sd}[F_j]\le d_j$. The leg
parameters are fitted on the legs of the development flights. Because the variance of one long leg can exceed that of
the two legs it replaces, removing a task can make a route infeasible under this model, so every route is filtered once
more before publication. With the cost-aware fallback, moment propagation accepts a given-up task if its nominal finish
meets the deadline and $\mathrm{E}[F_j]+z_j\,\mathrm{sd}[F_j]\le d_j$. Online calibration starts from one error model
fitted on both disturbance levels, with the weight of $n_0$ pairs, and refits Eq.~\eqref{eq:errmodel} at every call with
the pairs of the tasks completed so far, so it needs neither the disturbance level nor a level-specific safety factor.
The selected values are $c=0.45$ and 0.6, $z=1$ and 0.5 for moment propagation, and $n_0=50$ with $z=0.5$ for online
calibration, under mid and high disturbances. On development missions, the three static models differed by less than
0.2\,pp. CBBA bids the largest increase of the route value over all feasible insertion positions, plans from scratch at
every call and stops at convergence or at the hard budget. Without the bid warping of Johnson et
al.~\cite{johnson2012nonsubmodular}, it did not converge on some development snapshots. Within \mname{}, the stopping
rule uses a zero mean error, because the effective deadlines already contain it.

\subsection{Tuned competitors of studies 5 and 7}
In study~5, the proportional buffer was tuned with $c\in\{0.1,0.2,0.3,0.45,0.6,0.75,0.9,1.2\}$, with or without the
nominal fallback, and \mname{} without fallback with $z\in\{-0.5,0,0.5,1,1.5\}$, for every $\lambda\in\{0.5,1,2\}$ and
disturbance level. A grid that first ended at $c=0.3$ and $z=0$ was extended after both values were selected on its
edge, and the selection did not change. For $\lambda=1$, the selected buffer was \val{FSelPropOneMid} under mid and
\val{FSelPropOneHigh} under high disturbances, and \mname{} without fallback had \val{FSelCCnfOneMid} and
\val{FSelCCnfOneHigh}. Studies 5 and 6 ran the cost-aware fallback with the proportional cost model,
$\theta_j=\lambda/(1+\lambda)$, but scored the runs with $u_\lambda$, which charges every late arrival a fixed cost;
study~7a removes this mismatch. The fixed-cost rule of study~7a uses the mean reward $\bar w=2.1$ of the generator,
because the total reward of a mission is not known while tasks are still arriving, whereas $u_\lambda$ charges the
realised mean reward $R/n$ of the mission. Over the held-out missions, $R/n$ was \val{UCWnMean} on average, with a
coefficient of variation of \val{UCWnCv}\% and a range of \val{UCWnMin} to \val{UCWnMax}, so study~7a is also reported
with $u'_\lambda=U-\lambda\bar w\,n_{\mathrm{late}}/R$. In study~7a, the development selection of the tuned buffer
differed between $u_\lambda$ and $u^w_\lambda$ only at $\lambda=1$ under high disturbances, with $c=0.45$ for $u_1$ and
$c=0.6$ for $u^w_1$, both without fallback.

Study~7b keeps the error model, the safety factor and every tuned parameter of the calibrated high level. Speed and
service errors follow a Student-$t$ law with three degrees of freedom at the same standard deviation; rare long delays
add, with probability 0.1, an exponential delay with a mean of 90\,s to a leg, as a detour or a holding pattern would; a
persistent speed bias per UAV and mission with a standard deviation of 8\% correlates the errors along a route; and a
Pareto dissemination delay with the median of the lognormal one and a tail index of 1.5 replaces the lognormal delay.
It compares the nominal model, the constant and the proportional buffer, \mname{} with the nominal and with the
cost-aware fallback ($\lambda=1$, both cost models) and moment propagation with the cost-aware fallback.

\subsection{Protocol}
Studies 2, 3, 4, 5 and 6 used the held-out seeds 5000-5029, 7000-7029, 8000-8029, 10000-10029 and 12000-12029,
and studies 7a and 7b used 13000-13029 and 14000-14029. The tier C part of study~2 flew 30 missions with mid
disturbances and a latency factor of 50 (seeds 6000-6009), and study~4C used seeds 9000-9011. The protocols were not
lodged with an external registry. They are available with the code, together with the commit of every freeze and every
deviation. A tier C run in which a vehicle does not take off in time is repeated once, like a run with an invalid
real-time factor. Tier B runs four missions in parallel on an Intel Core i7-11800H laptop, with the methods of a
configuration run in turn. A separately written validator found no infeasible plan among \val{JValidatorPlans}
development plans. Because the stopping rules read a wall clock, the mean utility of a method varied by up to
\val{JRepeatDiffMax}\,pp between two runs of the same missions. Equivalence (H3) was tested with two one-sided tests.

\subsection{Further results of studies 2 to 4}
In study~2, the 90\% interval of the stopping rule minus a fixed 40\,ms was \val{KTostCI}\,pp, inside the
equivalence margin of $\pm0.5$\,pp. \mname{} without fallback served \val{KLateProp} tasks per mission after their
deadline, against \val{KLateForty} for ALNS with 40\,ms. In the tier C part of study~2, the Holm-adjusted $p$ of H4 was
\val{KCPForty}, and about \val{KCNPowerForty} missions per method would have been needed to confirm the difference that
tier B predicted. In study~4, the interval of \mname{} minus the proportional buffer was \val{XCIPropBufAll}\,pp
(Holm-adjusted $p={}$\val{QPPropBufAll}) and that of \mname{} minus moment propagation \val{XCIMPAll}\,pp. The loss of
moment propagation without fallback had an interval of \val{XCIFbMPAll}\,pp. \mname{} ran \val{QItMsCC} search iterations per millisecond and
moment propagation \val{QItMsMP}. Online calibration served \val{XDabsOnlCC}\,pp less than the model with a known
disturbance level, and the learned stopping rule added \val{XDabsFullCC}\,pp. In study~4C, the $p$-value of H14 was
\val{GPNom}. The advantage over the nominal model that sized the study was \val{XDabsTcDevNom}\,pp on development seeds.
Over the full tier B grid, the advantage in this regime was \val{XDabsFortyHighWHalf}\,pp, which would need about
\val{HNPowerNomGrid} missions per method to be confirmed in tier C. The interim analysis did not cross the
O'Brien-Fleming bound. The illustrative flights of Fig.~\ref{fig:demo} ran with the graphical interface, at a
real-time factor of \val{ZDemoRTFNom} to \val{ZDemoRTFProp}.

\subsection{Further results of studies 5 and 6}
In study~5, the advantage of the cost-aware fallback at $\lambda=1$ came from high disturbances, where it served
\val{XDabsFivePropOneHigh} and \val{XDabsFiveCCnfOneHigh}\,pp more than the tuned proportional buffer and the tuned
safety factor. Under mid disturbances, the methods were on par, and with the original windows both tuned competitors
were better, by \val{XDabsFivePropOneMidWOne} and \val{XDabsFiveCCnfOneMidWOne}\,pp. At $\lambda=0.5$, the interval
against the tuned buffer was \val{XCIFivePropHalf}\,pp. At $\lambda=2$, the cost-aware fallback stayed
\val{XDabsFiveCCfbTwo}\,pp above the nominal fallback and matched the tuned safety factor, while the tuned buffer was
better (\val{XCIFivePropTwo}\,pp) and moment propagation without fallback by \val{XDabsFiveMPnfTwo}\,pp. Tuned for
$\lambda=1$, the proportional buffer stayed \val{XDabsFivePropFixHalf}\,pp behind the cost-aware fallback at
$\lambda=0.5$. A wrong late cost also costs the cost-aware fallback, \val{XDabsFiveCAFixTwo}\,pp when it was run with
$\lambda=1$ and scored at $\lambda=2$, and when a late arrival costs nothing, the rule run with $\lambda=1$ served
\val{FUCostCCfb}\,pp less utility than the nominal fallback.

In study~6, CBBA on the nominal model reached almost the utility of ALNS (\val{WCICBBAFortyZero}\,pp), and \mname{}
served \val{WDabsCBBAEventZero}\,pp more than CBBA with the event trigger at $\lambda=0$. CBBA served \val{WLateCBBA}
tasks late per mission and \mname{} \val{WLateCA}. Under high disturbances with halved windows at $\lambda=0$, \mname{}
and CBBA were on par (\val{WDCBBAZeroHighWHalf}\,pp), because a fallback for $\lambda=1$ gives up tasks that a late
arrival would not penalise at $\lambda=0$. Within moment propagation, the cost-aware fallback lost
\val{WDabsMPFbZero}\,pp against the nominal fallback at $\lambda=0$ and gained \val{WDabsMPFbTwo}\,pp at $\lambda=2$.
The interval of the closed form minus moment propagation, both with the cost-aware fallback, was \val{WCICAMPCAOne}\,pp
at $\lambda=1$, and moment propagation ran \val{WItersMPCA} search iterations per call against \val{WItersCA}.

\subsection{Further results of study 7}
In study~7a, the fixed-cost rule scored \val{USCAFeeOneFee}\% with $u_1$, the proportional rule \val{USCAPropOneFee}\%
and the tuned buffer \val{USPropOneFee}\%; the Holm-adjusted $p$ of H24 was \val{UPHTwentyFour}. H26 had a
Holm-adjusted $p$ of \val{UPHTwentySix} for non-inferiority and $p={}$\val{UPSupHTwentySix} for superiority. The
fixed-cost rule scored \val{UDFeeCApropTwoFee}\,pp above the proportional rule with $u_2$, and the proportional rule
scored \val{UDabsFeeCApropOneW} and \val{UDabsFeeCApropTwoW}\,pp above the fixed-cost rule with $u^w_1$ and $u^w_2$.
With $u_2$, the fixed-cost rule minus the tuned buffer was \val{UDFeePropTwoFee}\,pp, whereas the proportional rule
stayed behind it (\val{UDPropPropTwoFee}\,pp). With $u^w_2$, the proportional rule stayed \val{UDabsPropPropTwoW}\,pp
behind the buffer tuned for that score. With $u'_1$, the differences of H24 and H25 were \val{UCDCApropOne} and
\val{UCDPropOne}\,pp. In study~7b, the nominal model served \val{UBUFortyGauss}\% under the calibrated errors and
between \val{UBUFortyOutlier}\% with rare long delays and \val{UBUFortyTthree}\% with Student-$t$ errors. The cost-aware
fallback scored \val{UDabsBCAPropTthree} to \val{UDabsBCAPropOutlier}\,pp more than the proportional buffer of study~4,
and rare long delays raised the mean finish-time error of \mname{} from \val{UBErrGauss} to \val{UBErrOutlier}\,s.


\section{Exact execution: stopping rules at equal computation}\label{sec:Sstatic}
\label{sec:static}
On the \val{NsnapTest} test snapshots, regret-2 repair finished within about one millisecond but left a regret of
\val{RepairRegret}\%. With a budget of 250\,ms, ALNS with the full budget reached \val{FixedRegret}\%, whereas
AnytimeSwarm reached \val{ASRegret}\% with \val{ASUsed}\% of the budget, that is, with \val{SavedPct}\% less computation.
Because the stopping rules use different amounts of computation, they are best compared at equal computation.
Table~\ref{tab:scale} makes this comparison for each fleet scale, with the fixed budget and the stagnation rule evaluated
at the mean computation that AnytimeSwarm used on the same scale. AnytimeSwarm reached a lower regret on every scale,
and all paired intervals lie below zero. The reduction was \val{ScMediumRedMin}\% to \val{ScMediumRedMax}\% on the
medium scale, \val{ScLargeRedMin}\% to \val{ScLargeRedMax}\% on the large scale and \val{ScXlargeRedMin}\% to
\val{ScXlargeRedMax}\% with 40 UAVs, a scale that was never used for training. Read in the other direction, the two
baselines needed between \val{ScMediumSpeedMin} and \val{ScSmallSpeedMax} times the computation of AnytimeSwarm to
reach its regret. Over all scales, AnytimeSwarm reached \val{EqRegretAS}\% against \val{EqRegretBaseMin}\% to
\val{EqRegretBaseMax}\% for four baselines that include a fixed budget and a stagnation patience tuned per fleet scale
on development data.

\inputtable{tabm_static_scale}

Fig.~\ref{fig:static} shows where this advantage comes from. On the medium and large scales, the quality-compute
frontier of AnytimeSwarm lies below those of fixed budgets and stagnation stopping over the range of practical
computation (Fig.~\ref{fig:static}a-c). Grouped by the improvement that a 500\,ms search can still find
(Fig.~\ref{fig:static}d), AnytimeSwarm spent \val{AllocASNone}\,ms on snapshots that cannot be improved and
\val{AllocASHigh}\,ms on those with the largest available improvement. Only \val{AllocASShare}\% of its computation went
to snapshots without any possible gain, compared with more than half for the baselines. The learned predictor is
essential for this allocation. It reached an area under the ROC curve~\cite{fawcett2006roc} of \val{AUCLearned} for the
event that the next 10\,ms bring an improvement, compared with \val{AUCOnline} for an extrapolation of the recent
improvement rate.

\begin{figure}[t]
\centering
\fig{fig_static_frontier}{\linewidth}\hspace{0.15in}\fig{fig_static_allocation}{\linewidth}
\caption{Exact execution on test seeds 2000-2029. (a)-(c) Quality-compute frontiers on identical search
trajectories. (d) Mean computation per snapshot, grouped by the improvement that a 500\,ms search can still
find.}\label{fig:static}
\end{figure}

\subsection{Reproduction on the evaluation platform}\label{sec:repro}
The exact-execution benchmark was produced on Windows. Before the execution-aware study, its development pipeline was
re-run on the platform of all later experiments (Ubuntu 24.04 under WSL2, identical Python and package versions). The
deterministic part was reproduced bit for bit, except for differences of one unit in the last place of some energy
values, which stem from the floating-point libraries of the two operating systems. The search ran
\val{JReproSpeedup}\% more iterations per millisecond on the new platform. The improvement predictor, re-trained on the
new timing traces, reached the same compute price by the rule fixed in the protocol, and the conclusions of the
equal-computation comparison agreed in sign and significance in \val{JReproMatched} comparisons. The frozen predictor
was used unchanged in all later experiments.


\section{Study 1: execution-aware stopping on the nominal planning model}\label{sec:Sstudy1}
Study~1 compares AnytimeSwarm-R with the stopping rules of Table~\ref{tab:methods} on the nominal planning model, on held-out seeds 3000-3029 in tier B (270 missions over three scales and three profiles, and 90
missions with 40 UAVs) and 4000-4009 in tier C (30 missions per method). Its primary matrix crosses two execution
levels (low, mid), two latency levels (none, high), two latency factors ($\kappa_B=10, 50$) and two triggers (event,
periodic every 60\,s). Its primary hypothesis is that AnytimeSwarm-R serves more reward than ALNS with 40\,ms,
AnytimeSwarm and the stagnation rule.
\subsection{Tier B primary matrix}\label{sec:tierb}
The primary matrix comprises \val{JNRunsMain} mission runs, that is, \val{JNMain} held-out missions in
\val{JNConfigsMain} configurations with eight methods. Table~\ref{tab:main} lists utility, lost tasks and computation per
method, and Fig.~\ref{fig:main} relates utility to computation and shows the difference in every configuration.

Search pays off clearly under uncertain execution. Regret-2 repair reached a utility of \val{JURepair}\%, and
AnytimeSwarm-R reached \val{JURisk}\%, which is \val{JDmRepair} percentage points (pp) or \val{JURelRepair}\% more. The
full budget of 250\,ms was the weakest search method, with \val{JUFixed}\%, and AnytimeSwarm-R served \val{JDFixed}\,pp
more reward than it. The loss of the long budget stems from late plans. A task counts as lost to delay when it is part of
the computed plan but can no longer be served on time when the plan is released. With $\kappa_B=50$, the full budget
lost \val{JLostKFiftyFixed} tasks per mission in this way, against \val{JLostKFiftyRisk} for AnytimeSwarm-R, that is,
\val{JLostRedFixed}\% fewer. Without its delay cost, AnytimeSwarm lost \val{JLostKFiftyNoDelay} tasks per mission in this
way, but its utility was not lower (\val{JUNoDelay}\% against \val{JUAS}\%). The delay cost thus prevents losses from
late plans without a measurable gain in utility on this matrix.

Among the remaining search methods, the differences are small. All seven search methods lie within
\val{JUSpreadSearch}\,pp of each other. AnytimeSwarm-R reached the highest mean utility with the least computation of the adaptive rules,
\val{JMsRisk}\,ms per call against \val{JMsAS}\,ms for AnytimeSwarm and \val{JMsStag}\,ms for the stagnation rule. The
reduction of \val{JMsRedAS}\% relative to AnytimeSwarm is significant (\val{JDMsAS}\,ms per call), so the second
hypothesis holds. The primary hypothesis does not hold at the level fixed in the protocol. The pooled advantage of
AnytimeSwarm-R was \val{JDFixedForty}\,pp over ALNS with 40\,ms, \val{JDAS}\,pp over AnytimeSwarm and \val{JDStag}\,pp
over the stagnation rule, with Holm-adjusted $p$-values of \val{JPFixedForty}, \val{JPAS} and \val{JPStag}. Only the
first interval excludes zero before correction. Per configuration and without correction for multiple testing,
AnytimeSwarm-R was better than ALNS with 40\,ms in \val{JWinFixedForty} of the \val{JNConfigsMain} configurations, better
than the stagnation rule in \val{JWinStag} and never different from AnytimeSwarm, which is consistent with chance
(Fig.~\ref{fig:main}b). The ablation that keeps only the latency term reached \val{JURiskZero}\%, close to AnytimeSwarm.
Together with the misspecification runs below, this suggests that a non-zero error spread mainly makes the rule stop
earlier, rather than that the calibration of the spread matters.

The fixed budget of 40\,ms is a strong reference, because it was matched to the mean computation of AnytimeSwarm on
development data. Its weakness lies in the allocation across fleet scales. AnytimeSwarm-R used \val{JMsScSmallRisk}\,ms
per call on the small scale and \val{JMsScLargeRisk}\,ms on the large scale. On the large scale, it served
\val{JDScLargeFixedForty}\,pp more reward than the budget of 40\,ms and was on a par with the full budget
(\val{JDScLargeFixed}\,pp). On the small scale, it served \val{JDScSmallFixed}\,pp more than the full budget and was on a
par with the budget of 40\,ms (\val{JDScSmallFixedForty}\,pp). A single fixed budget is therefore either too long for
small fleets or too short for large ones, whereas the adaptive rules find a suitable amount of computation on each scale
without tuning. The advantage over ALNS with 40\,ms was \val{JDFFKappa10}\,pp for $\kappa_B=10$ and
\val{JDFFKappa50}\,pp for $\kappa_B=50$, and \val{JDFFExecLow}\,pp under low and \val{JDFFExecMid}\,pp under mid
disturbances.

\inputtable{tabm_main}

\begin{figure}[t]
\centering
\fig{fig_main}{\linewidth}
\caption{Study 1, tier B. (a) Mean utility against mean planning time per call; regret-2 repair (\val{JURepair}\%, 1\,ms) is not shown. (b) AnytimeSwarm-R minus ALNS with 40\,ms and minus AnytimeSwarm per configuration, with 95\% intervals over \val{JNMain} missions.}\label{fig:main}
\end{figure}

\subsection{Replanning triggers}\label{sec:res-triggers}
The trigger mattered more than the stopping rule. Replanning every 60\,s in addition to the events raised the number of
planning calls from \val{JCallsEvent} to \val{JCallsPeriodic} per mission and the utility of every method, by
\val{JTrigMin}\,pp to \val{JTrigMax}\,pp, with all intervals above zero (\val{JTrigRisk}\,pp for AnytimeSwarm-R,
\val{JTrigFixedForty}\,pp for ALNS with 40\,ms). The third hypothesis therefore holds. Periodic replanning probably helps
because it corrects the accumulated deviations of the executed routes, which the event trigger only notices at the next
event. The advantage of AnytimeSwarm-R over ALNS with 40\,ms was \val{JDFixedFortyEvent}\,pp with the event trigger and
\val{JDFixedFortyPeriodic}\,pp with the periodic trigger. The two values are close, and the protocol did not test their
difference.

\subsection{Scalability to 40 UAVs}\label{sec:res-xl}
With 40 UAVs and 280 tasks, \val{JNXL} held-out missions were run under mid disturbances in \val{JXLNConf}
configurations of latency level and latency factor. All five methods reached between \val{JXLUAS}\% and
\val{JXLUStag}\%. AnytimeSwarm-R reached \val{JXLURisk}\% with \val{JXLMsRisk}\,ms per call. Its pooled difference was
\val{JXLDFixedOneEighty}\,pp to ALNS with 180\,ms, \val{JXLDAS}\,pp to AnytimeSwarm and \val{JXLDStag}\,pp to the
stagnation rule. After Holm correction over the four comparisons of the planned analysis, none of them is
significant ($p_{\mathrm{Holm}}=$ \val{JXLPAS} for AnytimeSwarm), so the fifth hypothesis is not supported either.
AnytimeSwarm-R again needed less computation than AnytimeSwarm (\val{JXLDMsAS}\,ms per call). The stagnation rule
stopped earliest, after \val{JXLMsStag}\,ms, and was significantly better than AnytimeSwarm-R in \val{JXLLossStag} of the
\val{JXLNConf} configurations, which had no latency and $\kappa_B=50$. On this scale, the rules used between \val{JXLMsStag} and 250\,ms
per call, and utility hardly changed within this range of computation. The equal-computation advantage that
AnytimeSwarm showed with 40 UAVs under exact execution (Table~\ref{tab:scale}) therefore did not translate into a
utility gain under uncertain execution.

\subsection{Sensitivity analyses}\label{sec:res-sens}
The sensitivity runs use the held-out missions of the primary matrix without latency and are descriptive. They vary the
execution level beyond the primary matrix, the latency factor and the trigger in \val{JNSensConf} configurations, and
the error model in \val{JNMisConf} further configurations.

Under high disturbances, which lie outside the range of the primary matrix but within that of the fitted error models,
utility fell to about \val{JSensUHighRisk}\% for all methods. The combination of high disturbances and a slow planner
($\kappa_B=50$) gave the largest advantage of the risk-aware cost in this study. AnytimeSwarm-R served
\val{JSensHighKFiftyAS}\,pp more reward than AnytimeSwarm and used \val{JSensMsHighKFiftyRisk}\,ms instead of
\val{JSensMsHighKFiftyAS}\,ms per call. Its difference to ALNS with 40\,ms was \val{JSensHighKFiftyFixedForty}\,pp and to
the stagnation rule \val{JSensHighKFiftyStag}\,pp, so the gain over AnytimeSwarm brought it to the level of the
baselines rather than beyond them. With $\kappa_B=1$, planning time hardly delays the plan, both delay costs are close to
zero because they scale with the shift $\kappa\tau$, and the two rules behave alike. AnytimeSwarm-R then differed from
AnytimeSwarm by \val{JSensKOneAS}\,pp and \val{JSensKOneMsDiff}\,ms per call.

Replanning when a UAV deviated by more than 30\,s from its planned finish time raised the number of calls from
\val{JSensCallsEvent} to \val{JSensCallsDev} per mission and the utility by \val{JSensTrigDevRisk}\,pp for
AnytimeSwarm-R and \val{JSensTrigDevFixedForty}\,pp for ALNS with 40\,ms. Like the periodic trigger, it helped both
stopping rules by a similar amount. Under the deviation trigger, the difference of AnytimeSwarm-R was
\val{JSensDevFixedForty}\,pp to ALNS with 40\,ms and \val{JSensDevAS}\,pp to AnytimeSwarm.

The utility of AnytimeSwarm-R was insensitive to the error model. With the error model of the low, mid or high level
under low and mid disturbances, its difference to AnytimeSwarm ranged between \val{JMisAllLow}\,pp and
\val{JMisAllMid}\,pp, the correctly specified model was not better than a wrong one, and the mean utilities of all
versions lay within \val{JMisSpread}\,pp. Setting the error model to zero gave
\val{JMisAllZero}\,pp. The four misspecification configurations also belong to the primary matrix, so the correctly
specified method was run twice on the same missions. The two runs differ only in the measured planning times, because
the stopping rule reads a wall clock. They gave identical utility in \val{JRepeatSameShare}\% of the missions, and the
mean utility of a method differed by up to \val{JRepeatDiffMax}\,pp between the runs. Differences of this size between
stopping rules should therefore not be read as systematic.

\subsection{Tier C: PX4 software-in-the-loop with Gazebo}\label{sec:res-tierc}
All \val{JNCRuns} held-out flights of tier C were completed, that is, \val{JNC} missions with four methods each. Two runs
did not reach their altitude within 90\,s and were repeated as specified. A further \val{JCNInvalid} runs had a real-time
factor between \val{JCRTFInvalidMin} and \val{JCRTFInvalidMax}, below the validity range of 1.6 to 2.2 stated in the
protocol. The protocol did not specify how to treat such runs, so they were repeated once with the same seed, like a
technical failure, and the original runs were kept. The repeated runs reached the same utility as the original ones, and
the comparisons below are identical with either set (\val{JCDFixedFortyOrig}\,pp and \val{JCDASOrig}\,pp with the
original runs). The real-time factor of the analysed runs was \val{JCRTF} on average and at least \val{JCRTFmin}. A
released plan reached all vehicles after \val{JCLat}\,s of simulation time on average, which is included in the expected
latency of AnytimeSwarm-R.

Table~\ref{tab:tierc} summarises the results. AnytimeSwarm-R reached a utility of \val{JCURisk}\%, against
\val{JCUFixedForty}\% for ALNS with 40\,ms, \val{JCUAS}\% for AnytimeSwarm and \val{JCUFixed}\% for the full budget.
The paired differences were \val{JCDFixedForty}\,pp to ALNS with 40\,ms and \val{JCDAS}\,pp to AnytimeSwarm, with
Holm-adjusted $p$-values of \val{JCPFixedForty} and \val{JCPAS}. The fourth hypothesis is therefore not supported. With
\val{JNC} missions, the intervals are much wider than in tier B and cannot resolve differences of the size seen there.
Two effects of tier B reappear in tier C. The full budget again had the lowest mean utility, although its interval
includes zero (\val{JCDFixed}\,pp to AnytimeSwarm-R), and it was the only method that lost a noticeable number of tasks
to late plans,
\val{JCDropFixed} per mission against at most \val{JCDropAS} for the others. AnytimeSwarm-R again computed least among
the adaptive rules, with \val{JCMsRisk}\,ms per call against \val{JCMsAS}\,ms for AnytimeSwarm, a reduction of
\val{JCMsRedAS}\% (\val{JCDMsAS}\,ms).

\inputtable{tabm_tierc}

The same \val{JNC} missions were also run in tier B with the rescaled parameters of tier C. Tier B predicted the outcome
of a flight well. The utility of a run in tier C correlated with that of the same run in tier B with $r=$ \val{JBCr}
(\val{JBCn} runs), and the mean utility per mission with $r=$ \val{JBCrMission} (Fig.~\ref{fig:bc}b). Tier C served
slightly more reward on average, \val{JBCUC}\% against \val{JBCUB}\%, so the kinematic model of tier B was, if anything,
slightly pessimistic for these missions. The small
differences between the methods were not predicted as well. The method means of the two tiers had a Kendall rank
correlation of \val{JBCTau}, and the sign of the per-mission difference between AnytimeSwarm-R and ALNS with 40\,ms
agreed in \val{JBCSignAgreeFixedForty}\% of the missions, counting ties as agreement only when both differences are zero.
Part of the per-run correlation is by construction, because both tiers draw the wind, speed and service disturbances
from the same random numbers, and much of it reflects the different difficulty of the missions. The correlation was
first computed with an erroneous pairing of runs in the analysis script, which was corrected after the results had
been seen; the hypothesis tests were not affected. In tier B, the advantage of AnytimeSwarm-R on these missions
was \val{JBgzDFixedForty}\,pp over ALNS with 40\,ms and \val{JBgzDAS}\,pp over AnytimeSwarm (Fig.~\ref{fig:bc}a). Both
tiers thus agree that the methods differ by at most a few percentage points on this scale, and that the full budget is
the weakest search method, while the order of the three short-computation methods is within the noise of both tiers.

\begin{figure}[t]
\centering
\fig{fig_bc}{\linewidth}
\caption{Study 1, tier B against tier C on the same \val{JNC} missions. (a) AnytimeSwarm-R minus each baseline, with 95\% intervals. (b) Utility of every run in tier C against the same run in tier B; the dashed line marks equality.}\label{fig:bc}
\end{figure}


\section{Study 2: regimes and software-in-the-loop details}\label{sec:Sstudy2}
Table~\ref{tab:s2main} lists the results of all methods and ablations of study~2, and Fig.~\ref{fig:Ss2main} shows their
paired differences together with the computation of every method.

\inputtable{tabm_s2_main}

\begin{figure}[t]
\centering
\fig{fig_s2_main}{\linewidth}
\caption{Study 2, tier B. (a) Utility of \mname{} without fallback minus each method, with 95\% intervals over \val{KNMissions} missions. (b) Mean utility against mean computation per call.}\label{fig:Ss2main}
\end{figure}

\subsection{Regimes}\label{sec:s2-regime}
Table~\ref{tab:s2regime} and Fig.~\ref{fig:s2regime} break the utility down by disturbance level, time-window factor and
latency factor. Per configuration and without correction for multiple testing, \mname{} without fallback was better than
ALNS with 40\,ms in \val{KWinForty} of the \val{KNConfigs} configurations and worse in \val{KLossForty}, better than the
constant buffer in \val{KWinBuf} and worse in \val{KLossBuf}, and better than the conventional planner in \val{KWinConv}
and worse in \val{KLossConv}. The advantage over ALNS with 40\,ms was \val{KDFortyKTen}\,pp, \val{KDFortyKFifty}\,pp and
\val{KDFortyKTwoHundred}\,pp for $\kappa_B=10$, 50 and 200, and the advantage over online prediction stayed between
\val{KDOnlKTen}\,pp and \val{KDOnlKTwoHundred}\,pp.

The advantage was largest with the original windows. Under high disturbances, \mname{} without fallback served
\val{KDFortyWOneHigh}\,pp more reward than ALNS with 40\,ms and \val{KDBufWOneHigh}\,pp more than the constant buffer.
The losses all occurred under high disturbances with halved windows. In this regime, \mname{} without fallback reached
\val{KUPropWHalfHigh}\%, \val{KDBufWHalfHigh}\,pp relative to the constant buffer and \val{KDFortyWHalfHigh}\,pp relative
to ALNS with 40\,ms. It served \val{KServedPropWHalfHigh} tasks per mission, of which \val{KLatePropWHalfHigh} late,
against \val{KServedBufWHalfHigh} and \val{KLateBufWHalfHigh} for the constant buffer. The chance-constrained model was
thus too cautious when the windows were short compared with the predicted finish-time error. The effective deadline of
many tasks then probably falls before the earliest possible arrival, so the planner gives up tasks that a less cautious
plan still serves in time. The value of $z$ was selected per disturbance level on development missions, averaged over
both window widths. Selecting it also per window width, or adapting it online, is a natural extension.

\inputtable{tabm_s2_regime}

\begin{figure}[t]
\centering
\fig{fig_s2_regime}{\linewidth}
\caption{Study 2: utility per regime against the latency factor, mean over \val{KNMissions} missions per point.}\label{fig:s2regime}
\end{figure}

\subsection{Tier C: PX4 software-in-the-loop with Gazebo}\label{sec:Ss2-tierc}
All \val{KCRuns} held-out flights of study~2 were completed, that is, \val{KCN} missions with four methods each. One run
failed to start and was repeated. As specified in the protocol, \val{KCInvalid} runs with a real-time factor outside
1.6 to 2.2 were repeated once with the same seed, and the original runs were kept. The laptop was also used for other
work during the campaign, and many repeats ran slower than the original runs. Of the \val{KCRuns} analysed runs,
\val{KCValid} lie in the validity range. The real-time factor was \val{KCRTF} on average, with quartiles of
\val{KCRTFq} and a minimum of \val{KCRTFmin}. Two sensitivity analyses that were not planned in the protocol therefore repeat the
comparisons with the original runs instead of the repeats, and with only those pairs of runs that both lie in the
validity range.

Table~\ref{tab:s2tierc} summarises the results. \mname{} without fallback reached a utility of \val{KCUProp}\%, against
\val{KCUForty}\% for ALNS with 40\,ms, \val{KCUBuf}\% for the constant buffer and \val{KCUConv}\% for the conventional
planner. The paired differences were \val{KCDForty}\,pp to ALNS with 40\,ms and \val{KCDConv}\,pp to the conventional
planner, with Holm-adjusted $p$-values of \val{KCPForty} and \val{KCPConv}. The fourth hypothesis is therefore not
supported. The secondary comparison with the constant buffer gave \val{KCDBuf}\,pp. The sensitivity analyses lead to
the same conclusion. With the original runs, the differences were \val{KCDFortyOrig}\,pp and \val{KCDConvOrig}\,pp.
With valid runs only, they were \val{KCDFortyValid}\,pp on \val{KCNFortyValid} missions and \val{KCDConvValid}\,pp on
\val{KCNConvValid} missions. No interval in the three analyses excludes zero.

\inputtable{tabm_s2_tierc}

The mechanism seen in tier B reappears in the counts of late and skipped tasks. \mname{} without fallback served
\val{KCLateProp} tasks late per mission, against \val{KCLateForty} for ALNS with 40\,ms, which is \val{KCLateRedForty}\%
fewer. It skipped no task at departure, whereas ALNS with 40\,ms skipped \val{KCSkipForty} and the conventional planner
\val{KCSkipConv} tasks per mission. It computed \val{KCMsProp}\,ms per call, against \val{KCMsForty}\,ms for the fixed
budget. The conventional planner called the search \val{KCCallsConv} times per mission instead of about
\val{KCCallsProp}, as in tier B.


The same \val{KCN} missions were also run in tier B with the rescaled parameters of tier C. As in study~1, tier B
predicted the utility of a flight closely, with a correlation of $r=$ \val{KBCr} over \val{KBCn} runs. On these
missions, however, even tier B did not resolve the differences between the methods (Fig.~\ref{fig:s2bc}). Its estimates
were \val{KBgzDForty}\,pp to ALNS with 40\,ms, \val{KBgzDBuf}\,pp to the constant buffer and \val{KBgzDConv}\,pp to the
conventional planner. The tier C intervals overlap them. The per-mission difference to ALNS with 40\,ms had a standard
deviation of \val{KCSdForty}\,pp in tier C. Confirming the tier B estimate of \val{KCEffForty}\,pp on these missions with
a power of 80\% would require about \val{KCNPowerForty} missions per method. Tier C thus shows that the proposed
replanner runs on an autopilot and a physics engine with the same computation and fewer late tasks, but it cannot
confirm the size of the gain measured in tier B.

\begin{figure}[t]
\centering
\fig{fig_s2_bc}{0.62\linewidth}
\caption{Study 2, tier B against tier C on the same \val{KCN} missions: \mname{} without fallback minus each baseline, with 95\% intervals.}
\label{fig:s2bc}
\end{figure}


\section{Study 3: regimes}\label{sec:Sstudy3}
Table~\ref{tab:s3} lists the results of study~3, and Fig.~\ref{fig:Ss3} shows its paired differences per regime.

\inputtable{tabm_s3}

\begin{figure}[t]
\centering
\fig{fig_s3}{\linewidth}
\caption{Study 3: utility of \mname{} with nominal fallback minus each method per regime, with 95\% intervals over \val{LNMissions} missions.}\label{fig:Ss3}
\end{figure}

\section{Study 4: regimes and reliability of the error model}\label{sec:Sstudy4}
Fig.~\ref{fig:replication} summarises the replication of the main comparisons on the held-out seeds of studies 2,
3 and 4.

\begin{figure}[t]
\centering
\fig{fig_replication}{\linewidth}
\caption{Replication on three independent sets of held-out seeds: utility of \mname{} minus each method, with 95\% intervals over 270 missions per study.}\label{fig:replication}
\end{figure}

Fig.~\ref{fig:s4} shows the paired differences of study~4 per regime. All three advantages of \mname{}, over
the constant buffer, the nominal model and the conventional planner, held in each of the four regimes.

\begin{figure}[t]
\centering
\fig{fig_s4}{0.85\linewidth}
\caption{Study 4: utility of \mname{} with nominal fallback minus each planning model per regime, with 95\% intervals over \val{QNMissions} missions.}\label{fig:s4}
\end{figure}

Fig.~\ref{fig:Scoverage} groups the tasks that were served in study~4 by the probability of meeting the deadline that
the error model of their disturbance level predicted at the last published plan, and compares it with the share served
on time. The plans of \mname{} and of the nominal model are evaluated with the same error model. Under mid
disturbances, the share served on time was \val{QCovObsCCMid}\% for \mname{}, against a mean predicted
probability of \val{QCovPredCCMid}\%. Under high disturbances, it was \val{QCovObsCCHigh}\% against
\val{QCovPredCCHigh}\%. Among the tasks with a predicted probability of at least one half, the shares were
\val{QCovObsAboveCCMid}\% and \val{QCovObsAboveCCHigh}\%, against predicted means of \val{QCovPredAboveCCMid}\% and
\val{QCovPredAboveCCHigh}\%. The model was thus slightly optimistic under mid and conservative under high disturbances,
and it ranked the tasks by their risk in both cases.

\begin{figure}[t]
\centering
\fig{fig_coverage}{\linewidth}
\caption{Study 4: share of served tasks on time against the predicted probability, for \mname{} with nominal fallback and the nominal model. The dotted line marks perfect calibration.}\label{fig:Scoverage}
\end{figure}

\section{Study 5: cost of a late arrival}\label{sec:Sstudy5}
Fig.~\ref{fig:Slatepen} scores the held-out runs of studies 3 and 4 again with a cost of $\lambda$ mean task rewards for
every late arrival, $u_\lambda=U-\lambda\,n_{\mathrm{late}}/n$. The analysis was done after study~4C, was not
planned in a protocol and used no new runs. The fallback of study~3 raised the number of tasks served late from
\val{PSThreeLateNoFb} to \val{PSThreeLateFb} per mission, and the fallback within moment propagation from
\val{PSFourLateNoFb} to \val{PSFourLateFb}. A comparison changes sign at the late cost
$\lambda^\ast=\Delta U/\Delta(n_{\mathrm{late}}/n)$. For the fallback, $\lambda^\ast$ was \val{PSThreeFbPropStar} in
study~3 and \val{PSFourMPMPnfStar} within moment propagation. In tier C, it was \val{PCFbNoFbStar}, with a difference of
\val{PCFbNoFbOne}\,pp at $\lambda=1$. The advantage over the constant buffer shrank with $\lambda$ and vanished at
\val{PSFourCCBufStar}, because the buffer of 60\,s served fewer tasks late than \mname{} with nominal fallback (\val{QLateBuf} against \val{QLateCC} per mission). Against the proportional buffer, the
proposed model with fallback stayed behind at every $\lambda$, by \val{PSFourCCPropBufOneAbs}\,pp at $\lambda=1$
(\val{PSFourCCPropBufOne}). Study~5 (Section~\ref{sec:study5} of the main paper) was designed from these results.

\begin{figure}[t]
\centering
\fig{fig_latepen}{\linewidth}
\caption{Studies 3 and 4 re-scored with a late cost of $\lambda$ mean task rewards: score differences with 95\% intervals over 270 missions (exploratory, no new runs).}\label{fig:Slatepen}
\end{figure}

\section{Calibration of the predicted probability}\label{sec:Scalib}
Every tier B run since study~4 bins its executed tasks by the probability of meeting the deadline that the error model
of the disturbance level predicted at the last published plan, and counts how many of them met it. From these bins,
Fig.~\ref{fig:Scalib} shows the reliability of the prediction for \mname{} in studies 4-6 and, in its right
panel, under the shifted execution errors of study~7b. Over studies 4-6, the expected calibration error,
$\mathrm{ECE}=\sum_i|k_i-p_i|/N$ with $k_i$ the tasks on time and $p_i$ the sum of the predicted probabilities in
decile $i$, ranged from \val{VCalEceMin} to \val{VCalEceMax}\,pp. The upper tail, a predicted probability of at least
0.8, is where a chance constraint acts. There, \val{VCalObsHiFiveMid}\% of the tasks were served on time under mid
disturbances in study~5, against a predicted \val{VCalPredHiFiveMid}\%, and \val{VCalObsHiFiveHigh}\% under high
disturbances, against \val{VCalPredHiFiveHigh}\%. Below the tail, the model was conservative under high disturbances,
where it predicted a lower probability than was observed, and slightly optimistic under mid disturbances. Under the
shifted errors of study~7b, the ECE was \val{VCalEceSGauss}\,pp for the calibrated Gaussian errors,
\val{VCalEceSTthree}\,pp for Student-$t$ errors, \val{VCalEceSOutlier}\,pp with rare long delays,
\val{VCalEceSBias}\,pp with a speed bias per UAV and \val{VCalEceSPareto}\,pp with Pareto latency, with Brier scores of
\val{VCalBrierSGauss}, \val{VCalBrierSTthree}, \val{VCalBrierSOutlier}, \val{VCalBrierSBias} and
\val{VCalBrierSPareto}. In the upper tail, the share on time under these conditions was \val{VCalObsHiSGauss},
\val{VCalObsHiSTthree}, \val{VCalObsHiSOutlier}, \val{VCalObsHiSBias} and \val{VCalObsHiSPareto}\%, against predicted
means of \val{VCalPredHiSGauss}, \val{VCalPredHiSTthree}, \val{VCalPredHiSOutlier}, \val{VCalPredHiSBias} and
\val{VCalPredHiSPareto}\%.

\begin{figure*}[t]
\centering
\fig{fig_calib}{\linewidth}
\caption{Reliability of the predicted probability of meeting the deadline for \mname{}. (a) Mid and (b) high disturbances in studies 4-6. (c) Study 7b with the calibrated errors and four departures from the error model. The dotted line marks perfect calibration.}\label{fig:Scalib}
\end{figure*}

\section{Sensitivity to the resampling unit}\label{sec:Scluster}
The primary tests resample missions. The nine missions of one seed have independent scenarios but share the execution
disturbances, which are keyed to the seed, the UAV and the task. Table~\ref{tab:Scluster} repeats every tier B primary
test of studies 2-7 with a paired bootstrap over the \val{VClNClusters} seeds. The intervals were \val{VClRatioMin} to
\val{VClRatioMax} times as wide as over missions (median \val{VClRatioMed}). They were mostly narrower, because every
seed contains one mission of each fleet scale and event profile, so that resampling seeds removes the variation between
these strata. \val{VClNChanged} of the \val{VClNTests} outcomes changed. In H9, the advantage of the proportional buffer
over \mname{}, which had missed significance over missions, became significant over seeds (interval
\val{VClHNineCI}\,pp for \mname{} minus the buffer, Holm-adjusted $p={}$\val{VClHNineP}). The tier C
hypotheses have one mission per seed and are not affected.

\inputtable{tabm_cluster}

\section{Study 7: further results}\label{sec:Sstudy7}
% Study 7 further results (docs/16_PROTOCOL_V7.md; numbers from experiments/analyze_study7.py, macros \val{U...}).
Study~7a ran the cost-aware fallback with both cost models for $\lambda\in\{0.5,1,2\}$ on seeds 13000-13029. The
competitors were selected on development missions per score, late cost and disturbance level. For $u_\lambda$, the
proportional buffer was \val{USelPropHalfFee}, \val{USelPropOneFee} and \val{USelPropTwoFee} for $\lambda=0.5$, 1 and
2 (mid / high disturbances), the same selection as in study~5. For $u^w_\lambda$, it was \val{USelPropHalfW},
\val{USelPropOneW} and \val{USelPropTwoW}. Table~\ref{tab:Ss7a} gives the scores of the main methods.

\begin{table}[t]
\caption{Study 7a: mean score (\%) of the two rules and of the tuned proportional buffer.}
\label{tab:Ss7a}
\centering\footnotesize
\setlength{\tabcolsep}{3pt}
\begin{tabular}{@{}lcccccc@{}}
\toprule
& \multicolumn{3}{c}{$u_\lambda$ (fixed cost)} & \multicolumn{3}{c}{$u^w_\lambda$ (cost $\propto w_j$)} \\
\cmidrule(lr){2-4}\cmidrule(lr){5-7}
Method & $\lambda=0.5$ & 1 & 2 & 0.5 & 1 & 2 \\
\midrule
Fixed-cost rule & \val{USCAFeeHalfFee} & \val{USCAFeeOneFee} & \val{USCAFeeTwoFee} & \val{USCAFeeHalfW} &
\val{USCAFeeOneW} & \val{USCAFeeTwoW} \\
Proportional rule & \val{USCAPropHalfFee} & \val{USCAPropOneFee} & \val{USCAPropTwoFee} & \val{USCAPropHalfW} &
\val{USCAPropOneW} & \val{USCAPropTwoW} \\
Tuned proportional buffer & \val{USPropHalfFee} & \val{USPropOneFee} & \val{USPropTwoFee} & \val{USPropHalfW} &
\val{USPropOneW} & \val{USPropTwoW} \\
\bottomrule
\end{tabular}
\end{table}

Against the nominal fallback, the fixed-cost rule scored \val{UDFeeCCfbOneFee}\,pp more with $u_1$ and
\val{UDFeeCCfbTwoFee}\,pp more with $u_2$; the proportional rule scored \val{UDPropCCfbOneW}\,pp more with $u^w_1$.
By regime at $\lambda=1$, the fixed-cost rule minus the tuned buffer with $u_1$ was \val{UXRFeeMidWOne}\,pp under mid
disturbances with the original windows, \val{UXRFeeMidWHalf}\,pp with halved windows, and \val{UXRFeeHighWOne} and
\val{UXRFeeHighWHalf}\,pp under high disturbances with the original and the halved windows. As in study~5, the
advantage came from high disturbances and short windows, and with mid disturbances and the original windows the tuned
buffer was ahead.

The fixed-cost rule assumes a cost of $\lambda\bar w$ per late arrival with the mean reward $\bar w=2.1$ of the
generator, whereas $u_\lambda$ charges $\lambda R/n$ with the realised mean reward of the mission, which ranged from
\val{UCWnMin} to \val{UCWnMax} over the held-out missions. Table~\ref{tab:Ss7const} repeats the comparisons of the
fixed-cost rule with $u'_\lambda=U-\lambda\bar w\,n_{\mathrm{late}}/R$, which charges exactly $\lambda\bar w$. All
differences changed by at most a few hundredths of a point, and no conclusion changed. Using the realised $R/n$ in the
rule itself is not possible during a mission, because $R$ includes the tasks that have not yet arrived.

\inputtable{tabm_s7const}

Study~7b ran on seeds 14000-14029 under high disturbances. Table~\ref{tab:s7} gives the utility and the score $u_1$ of
every method under the calibrated Gaussian errors and the four departures. The mean finish-time error of the proposed
model with the cost-aware fallback was \val{UBErrGauss}, \val{UBErrTthree}, \val{UBErrOutlier}, \val{UBErrBias} and
\val{UBErrPareto}\,s, and its 95th percentile \val{UBErrPGauss}, \val{UBErrPTthree}, \val{UBErrPOutlier},
\val{UBErrPBias} and \val{UBErrPPareto}\,s. On $u_1$, the fixed-cost rule scored \val{UDabsBCACAfeeTthree} to
\val{UDabsBCACAfeeGauss}\,pp above the proportional rule, and moment propagation with the cost-aware fallback
\val{UDabsBCAMPCAGauss} to \val{UDabsBCAMPCATthree}\,pp above the closed form, with every interval including zero in
both cases. Moment propagation served less utility than the closed form under every condition.

\inputtable{tabm_s7}


% Reference list generated with BibTeX (style elsarticle-num, references_ast.bib with LTWA journal
% abbreviations) and pasted here; regenerate with the scripts in _build/.
\begin{thebibliography}{10}
\expandafter\ifx\csname url\endcsname\relax
  \def\url#1{\texttt{#1}}\fi
\expandafter\ifx\csname urlprefix\endcsname\relax\def\urlprefix{URL }\fi
\expandafter\ifx\csname href\endcsname\relax
  \def\href#1#2{#2} \def\path#1{#1}\fi

\bibitem{mccullagh1989glm}
P.~McCullagh, J.~A. Nelder, Generalized Linear Models, 2nd Edition, Chapman and
  Hall, London, 1989.
\newblock \href {https://doi.org/10.1007/978-1-4899-3242-6}
  {\path{doi:10.1007/978-1-4899-3242-6}}.

\bibitem{savelsbergh1992vrptw}
M.~W.~P. Savelsbergh, The vehicle routing problem with time windows: Minimizing
  route duration, ORSA J. Comput. 4~(2) (1992) 146-154.
\newblock \href {https://doi.org/10.1287/ijoc.4.2.146}
  {\path{doi:10.1287/ijoc.4.2.146}}.

\bibitem{ortools}
L.~Perron, V.~Furnon, {OR-Tools},
  \url{https://developers.google.com/optimization}, version 9.15 (2026).

\bibitem{voudouris1999gls}
C.~Voudouris, E.~Tsang, Guided local search and its application to the
  traveling salesman problem, Eur. J. Oper. Res. 113~(2) (1999) 469-499.
\newblock \href {https://doi.org/10.1016/S0377-2217(98)00099-X}
  {\path{doi:10.1016/S0377-2217(98)00099-X}}.

\bibitem{svegliato2018meta}
J.~Svegliato, K.~H. Wray, S.~Zilberstein, Meta-level control of anytime
  algorithms with online performance prediction, in: Proceedings of the
  Twenty-Seventh International Joint Conference on Artificial Intelligence
  (IJCAI-18), IJCAI Organization, 2018, pp. 1499-1505.
\newblock \href {https://doi.org/10.24963/ijcai.2018/208}
  {\path{doi:10.24963/ijcai.2018/208}}.

\bibitem{clark1961greatest}
C.~E. Clark, The greatest of a finite set of random variables, Oper. Res. 9~(2)
  (1961) 145-162.
\newblock \href {https://doi.org/10.1287/opre.9.2.145}
  {\path{doi:10.1287/opre.9.2.145}}.

\bibitem{johnson2012nonsubmodular}
L.~B. Johnson, H.-L. Choi, S.~S. Ponda, J.~P. How, Allowing non-submodular
  score functions in distributed task allocation, in: Proceedings of the 51st
  IEEE Conference on Decision and Control (CDC), 2012, pp. 4702-4708.
\newblock \href {https://doi.org/10.1109/CDC.2012.6425867}
  {\path{doi:10.1109/CDC.2012.6425867}}.

\bibitem{fawcett2006roc}
T.~Fawcett, An introduction to {ROC} analysis, Pattern Recognit. Lett. 27~(8)
  (2006) 861-874.
\newblock \href {https://doi.org/10.1016/j.patrec.2005.10.010}
  {\path{doi:10.1016/j.patrec.2005.10.010}}.

\bibitem{ropke2006adaptive}
S.~Ropke, D.~Pisinger, An adaptive large neighborhood search heuristic for the
  pickup and delivery problem with time windows, Transp. Sci. 40~(4) (2006)
  455-472.
\newblock \href {https://doi.org/10.1287/trsc.1050.0135}
  {\path{doi:10.1287/trsc.1050.0135}}.

\bibitem{boddy1994deliberation}
M.~Boddy, T.~L. Dean, Deliberation scheduling for problem solving in
  time-constrained environments, Artif. Intell. 67~(2) (1994) 245-285.
\newblock \href {https://doi.org/10.1016/0004-3702(94)90054-X}
  {\path{doi:10.1016/0004-3702(94)90054-X}}.

\bibitem{zilberstein1996using}
S.~Zilberstein, Using anytime algorithms in intelligent systems, AI Mag. 17~(3)
  (1996) 73-83.
\newblock \href {https://doi.org/10.1609/aimag.v17i3.1232}
  {\path{doi:10.1609/aimag.v17i3.1232}}.

\bibitem{jabali2015selfimposed}
O.~Jabali, R.~Leus, T.~Van~Woensel, T.~de~Kok, Self-imposed time windows in
  vehicle routing problems, OR Spectrum 37~(2) (2015) 331-352.
\newblock \href {https://doi.org/10.1007/s00291-013-0348-1}
  {\path{doi:10.1007/s00291-013-0348-1}}.

\end{thebibliography}
\end{document}
